\chapter{Summary}

This projects goal is to create a game engine, similar in design to currently available big commercial engines.
As this is a large undertaking that far exceeds the limits and practicality of a thesis project, I will be focusing on giving the project a good foundation. This means considerable amount of time will be spent investigating solutions to issues other projects would not care about. By the end I hope to have a solid foundation and core systems for the engine that I plan to use and extend in the coming years.
Planned parts:

\begin{itemize}
    \item CI/CD pipelines
    \item Choosing a build system that is usable with large projects
    \item Dynamic DLL loading of engine modules
    \item Basic Entity system
    \item Reflection system for runtime code metadata
\end{itemize}

The project is a collaboration between me and a friend of mine, she will be writing the renderer and donating some other bits of code. The parts that I will be presenting are entirely my own.

The project is hosted on github and is MIT licensed.


\chapter{Intro}

Most well known formas are mobile games like: Candy Crush or bigger budget console games like "The last of us" \footnote{This game series recently even got a really successful TV adaptation by HBO}
These pieces of software need a lot of different disciplines to come together in order to create them. At the bare minimum, one needs to be a good software engineer to create the code for the game and also ba able to design a good and engaging game. But when it comes to large production games, so-called "AAA" or "triple A" games the number of necessary tasks increases. For these games on top of the programming and game design new tasks come in like story writing, 3D modelling, environment design, level design, music, and countless others. For these games, there is an expectation to excel in every category or even push what is possible.

Because video games as so complex, there is a need to share some of the code between projects, as starting from scratch for each is impractical. This is possible because many of the functionality needed to create a game is the same for every game. For example, almost all games need a way to show 3D models on the display, and they need to process input from the player, regardless of what the actual gameplay is. These more abstract, not game specific functionalities, can be shared between projects. 

These sharable functionalities are what form many of the game development code libraries. These are not runnable code by themselves and are only for the programmers to help and accelerate creating a game. Game engines are even more complex and usually include editors and tools to create the assets that are going to be used in the game. This usually includes a world editor to create the 2D or 3D environment the player can explore. Game engines are responsible for running the actual game application.


\section{What is a game engine?}
A game engine is a highly complex software or library that allows the creation of computer games. They can vary is what functionality they provide and how the game's creator interacts with them. They can have graphical interfaces and easy to understand UI, or be code only and need programming knowledge to use. The most wildly known games engines (like Unity and UE4) provide a single editor application that the user (the game maker) interacts with. These programs have many built in tools that help with creating a game, like level editor tools, material editors, some even provide node based scripting tools.

There are game engines that are available for free and used by many companies to building games. While others use either a purpose built engines for their games or one that the company developed in-house.


\section{How games work}
Games are real time graphical applications, meaning that they redraw the contents of the screen fast enough where they form a continuous animation. Films work the same way but as they are not interactive they can work with 25 frames per second (FPS), but as games need to react to the user's input they should strive to draw as many frames as the hardware can process. For games developed for PCs the recommended minimum is 60 frames per second. This means the game has to finish redrawing and updating everything in the game world 60 times a second.

Engines are the core of a games code, they provide many of the utility functions that are needed to draw the game, gather input, load files and many more. Just as there are many different styles of games, there are many different game engines. Some are purpose built engines for a specific game with no intention to be reusable. While others are general purpose engines that can be used to make a wide range of games. These game engines will also vary in what they provide, for example an engine for a single player game will probably not have any networking code.

In many cases companies will have their own engine developed for their own games, this is useful for them as they can modify the core engine code for each game as needed. This can result in higher integration and optimization between the engine and the game. Not many of these so called "in house" engines are available to the public to see. The biggest two publicly available that fit this description are Valve's Source engine and Epic Gemes' Unreal Engine, both of these engines have a long history going back to the 90's and the start of the computer gaming industry.

Because there are countless engines most of witch is not available online I will be only using a few of the most known engines for comparison.
The 3 most popular freely available game engines:
- Unity
- UE4/5
- Godot
On top of these I will also include Valve's Source engine as it is source available and in some cases has simpler implementations of certain tools/functionality

\subsection{Unity}

Unity is a fully multi purpose engine capable of 2D and 3D graphics with physics, and modern rendering features. It's biggest selling point is the large amount of platforms it supports, games made with it can run on any modern PC and every console.

I personally have the most experience with it out of the 3, and it is one of the reasons why I started this project. While using unity, I often found my self wanting to do more complex things than what the engine allowed. This meant that I ended up fighting with the engine more than just using it. This became extremely apparent during my last laboratory project (Ön lab, Unity event system).

\subsection{Unreal 4, 5}

Unreal is one of the oldest engines still under development. Its biggest feature is the high level of visual quality it can provide. Unreal is mostly focused on the AAA style market and FPS / TPS games. It's other big selling point is that it is source available, this means that if you want you can edit the source code and compile a version for your own use (while adhering to the license).

Uneral's problem for my is that it is highly focused on multiplayer FPS games, and some of the settings and facilities it provides make no sense outside of that use case. On top of that the editor and development environment is a large install size witch makes it difficult to use on slower devices.

\subsection{Godot}

Godot is relatively the newest game engine out there, it is developed fully in the open and is MIT licensed. It has full support of 3d graphics and is constantly improving the visual fidelity it can provide. The fact that it is developed fully in the open gives it an advantage when it comes to extensibility, but also makes the development harder to fund.

It has C\#, Cpp, and their own GScript options for writing your game code this makes developing complex system, and simpler game logic both easy.

\subsection{Source}

Source is Valve's game engine and was used in almost all of their games, like Half-Life 2, portal1/2, and many more.

It follows and older style of game engine design where the runtime of the engine is almost completely separate from the editing tools. For example it has a level editor that uses a separate way to draw the designs, and needs to export it in a format that the runtime can understand.

In my experience it has the most efficient level editor, but also the one that feels the oldest. The engine has a new version used in newer games made by Valve. This Source 2 engine still shares a lot of aspects with the old version while providing a much more integrated development experience. Sadly this version is not yet source available and is hard to get access to in order to develop games with.


\section{The project}

Base idea

\subsection{My history with game programming}
I have been writing games engines and game related software for a long time.
My first time using a professional game engine was back in 2015 with the 4.5-5.0 versions of Unity, this was at the start of my 6th grade in elementary school. From there, I mainly used unity for a staggering amount of side projects that never actually got finished or were mostly library style projects. \ednote{Should maybe cut this down?}

In the last year of high school, I have started working on my first game engine in C\#, called Cryogenic engine.
The goal of this project was to learn how to use the then upcoming Vulkan graphics library.
The project never got that far as I have not fully understood how big of a task implementing a Vulkan based renderer is.
This project got put on hold as I started my university studies.

In university's first semester, I have created a small game engine for my ANSI C 101 course as an extra assignment.
This project introduced me to how game engines in lower level languages work and what problems they need to solve.
This engine was called Relic engine ( Joke on C being old :) ),
it was a relatively simple 2D engine that was able to load in tile based map from the widely used 3rd party external editor called Tiled \cite{tiled}

In the next semester in my C++ 101 course, I have started writing a game engine for that course's extra assignment as well. This first started with the idea of "Let's make Relic engine in C++" [Addendum 2d shadows].
With C++ I've gotten access to classes and with this more structure for my code.
This lead to a massive feature creep, witch made the it so that only the input system got finished by the end of the semester. But I continued to chip away on it after that. This became the old ShadowEngine.

When I have started a DirectX12 programming course I have began creating a dx12 based renderer for this old version of the engine. This while worked, was a really simplistic implementation and plagued me for the next few years.

\subsection{The new engine}

From all of these older attempts i have learned a lot of important lessons. The most important two things are that writing a renderer is hard. They are one of the most complex parts of a game engine and need to handle a lot of different data. They are also hard to know if you have written them right. You might only come across a use case that the design can't handel after it is too late and needs major rewrites. 

The second major realization was that making the foundations of the engine more thought out is really important. In the first iteration of this engine I have jumped many steps in order to get to a more "fun" part of the development. This has made some core components rushed and hard to use, while other unrelated sections got tangled up.

This is when, while talking with friends, the idea of a rewrite came up. As my thesis project was coming up it seamed like a great opportunity to start the project. At this point I was looking around to see what kind of existing rendering solutions could I use. There are a number of libraries out there that provide varying levels and types of rendering. I have decided to accept my friends offer to make a renderer for me while deciding to create the rest of the engine in a way to allow for easy adaption of any existing solution. 

The requirements
\begin{itemize}
    \item Modules
    \item Inter module communication
    \item Entities
    \item Editor runtime setup
\end{itemize}

\subsection{Licensing}

The project is currently licensed under MIT and will stay that way for the foreseeable future
